\section{Function-Space Overlap: Neural Operators and Mergeable Sketches}\label{app:operator-overlap}

Section~\ref{sec:fno-primer} introduced neural operators as reusable
function-space maps, and Section~\ref{sec:mergeable-sketches} introduced
mergeable sketches as encode/merge/decode summaries. This appendix
records the exact overlap. The point is not that every neural operator
is a mergeable sketch, or that every classical sketch is implemented by
a transformer. The point is that the paper uses named subspaces and
named intersections.
Appendix~\ref{app:kovachki-finite-dimensionalization} records the
separate Kovachki Lemma~21/22 finite-dimensionalization route; this
appendix uses that route only through the resulting approximation and
local-law subspace interfaces.

\subsection{Notation Map}

Table~\ref{tab:ctreepo-neural-operator-notation} adapts the notation
dictionary role of \citet[Table~10]{Kovachki23NeuralOperator} to the
objects used in this paper. The key shift is that our ``domain'' is not a
PDE mesh alone: it is the index set of the realized call currently being
processed, whether that call is a raw leaf, a pair of child summaries, or
an on-range summary being re-used.

\begin{table}[t]
\centering
\footnotesize
\setlength{\tabcolsep}{3pt}
\begin{tabular}{p{0.18\linewidth}p{0.34\linewidth}p{0.40\linewidth}}
\toprule
Neural-operator notation & Meaning in C-TreePO & Examples in this paper \\
\midrule
$D$, $x\in D$ &
Index domain and locations inside one operator call. &
Markov token positions; HLL stream/register coordinates; manifesto character or quasi-sentence offsets. \\
$a\in\mathcal{A}$ &
Input function or state presented to the operator. &
A leaf block, concatenated child summaries, or an on-range summary; Markov observed-token sequence; HLL stream/register state; RILE span plus rubric prompt. \\
$u\in\mathcal{U}$ &
Target output function, summary state, or readout-facing state. &
Markov $(\text{count},\text{first},\text{last})$ state or root count; HLL register vector or estimate; RILE evidence summary or score vector. \\
$D_j$ &
Discretization/refinement at which the same operator family is evaluated. &
Chunk size and leaf geometry; tokenizer/grid length for FNO or transformer calls; HLL precision as the register-state resolution. \\
$\mathcal{G}^{\dagger}:\mathcal{A}\to\mathcal{U}$ &
Ideal target operator before approximation. &
Exact theorem-backed summarizer; hand-designed mergeable sketch; oracle-preserving Markov boundary-state merger. \\
$\mathcal{G}_{\theta}$, $g$ &
Realized learned operator used by the tree. &
1D FNO leaf/merge network; learned HLL-style merge/readout; prompted LLM leaf and merge operator. \\
$\mu$ &
Distribution over inputs or realized calls. &
Synthetic Markov DGP; HLL benchmark stream distribution; Manifesto corpus and sampled audit-node distribution. \\
$v_t$, $d_v$, $P,Q,K,W,\sigma$ &
Hidden representation dimension and equation-(6) layer components. &
FNO latent state and Fourier kernel; transformer attention/kernel layer plus feed-forward maps; MLP merge/readout adapters. \\
$K_{\mathrm{call}}$ &
Compact realized-call set on which approximation is required. &
The finite/compact collection of leaf calls, merge calls, and on-range re-summary calls generated by one audited tree run. \\
\bottomrule
\end{tabular}
\caption{C-TreePO notation map for the neural-operator layer, paralleling the operator-learning notation dictionary of \citet[Table~10]{Kovachki23NeuralOperator} while making the paper's three examples explicit.}
\label{tab:ctreepo-neural-operator-notation}
\end{table}

\subsection{The Diagram}

Fix an ambient operator space
\[
\mathrm{Op}(\mathcal{A},\mathcal{U})
  := \{\,\mathcal{G}:\mathcal{A}\to\mathcal{U}\,\}.
\]
The formalized diagram is:
\[
\begin{tikzcd}[column sep=large,row sep=large]
\text{finite transformer encoders}
  \arrow[r, hook]
  & \mathrm{Eq6NO}
  \arrow[r, hook]
  & \mathrm{NO}
  \arrow[d, "\cap\mathrm{MS}" description] \\
\text{mergeable sketches}
  \arrow[r, hook]
  & \mathrm{ExactLL}(T,\fstar)
  \arrow[r, hook]
  & \mathrm{ApproxLL}(T,\fstar;\varepsilon)
\end{tikzcd}
\]
Here $\mathrm{Eq6NO}$ is the equation-(6) class from
Equation~\eqref{eq:equation6-no}, $\mathrm{MS}$ is the class of
operators induced by mergeable sketches with exact sketch-local
witnesses, $\mathrm{ExactLL}$ is the exact C1/C2/C3 subspace, and
$\mathrm{ApproxLL}$ is its budget-relaxed version. The right vertical
label is the certified overlap:
\[
\mathrm{MNO}_{\mathcal{N},\fstar}
  :=
  \mathcal{N}\cap\mathrm{MS},
\]
where $\mathcal{N}$ is the chosen neural-operator class.

\subsection{Transformer Inclusion}

The transformer inclusion is architecture-level. Following Kovachki
Proposition~6, a pre-normalized single-head attention block agrees with
a discretized kernel layer on the block's active coordinates; stacking
attention, residual/local maps, and pointwise feed-forward layers
produces an equation-(6) neural operator. The membership statements we
use are therefore: (i) single-head attention equals a discretized kernel
layer, (ii) a transformer block is a neural-operator layer, (iii) a
finite transformer encoder stack is an equation-(6) neural operator,
and (iv) the encoder stack belongs to the equation-(6) class. These
membership claims operate at the architecture level only; they do not
certify tokenizers, runtime serving details, or a specific model
checkpoint.

\subsection{Mergeable-Sketch Inclusion}

The classical lane is broader than HLL. A sketch is a triple
(encode, merge, decode); the sketch-local-law status is established
separately for HLL, Count-Min, KLL quantiles, GK's sequential quantile
surface, bigram sketches, and Markov-count sketches. HLL is the
empirical parity anchor because it gives a clean byte-identity test,
but the theorem surface is the general mergeable-sketch class.

The key class is $\mathrm{MS}$: a realized operator $g$ belongs to it
when $g$ is the decode-after-encode action of a sketch whose encoder
preserves oracle value at leaves, whose merge preserves it across
siblings, and whose decoder is compatible with the oracle. Any such
sketch-induced operator lies in the exact local-law subspace on every
tree.

\subsection{Certified Overlap}

For a chosen neural-operator class $\mathcal{N}$, the certified overlap is
\[
\mathrm{MNO}_{\mathcal{N},\fstar}
  \;:=\;
  \mathcal{N} \cap \mathrm{MS}(\fstar),
\]
and this overlap sits inside the exact local-law neural-operator subspace
$\mathrm{ExactLL}_{\mathcal{N}}(T,\fstar)$. Separately, the older
Proposition~\ref{prop:mergeable-reduction} route yields the same overlap
through a different argument: any deterministic operator carrying
A1/A2/A3 witnesses becomes a classical mergeable summary, and can then be
represented by a sketch adapter whose encoder is the operator itself,
whose merge is the operator applied to the concatenation of child
summaries, and whose decoder is the identity.

\subsection{Oracle \texorpdfstring{$+$}{+} Projection Decomposition}\label{app:oracle-plus-projection}

The operational learned setting usually lands in the approximate, not
exact, subspace. The formal objective therefore separates two roles:
oracle fit and projection toward local-law validity. The main-text
equation~\eqref{eq:oracle-projection-objective} writes the objective
with a single scalar $\lambda\in[0,1]$, weighting the three audited
budgets uniformly. A more general per-law-weighted form with
non-negative shares $\rho_1,\rho_2,\rho_3$ reads
\begin{equation}\label{eq:oracle-projection-objective-generalized}
J_\lambda^{\boldsymbol{\rho}}(g)
  =
  (1-\lambda)\,L_{\mathrm{oracle}}(g)
  +
  \lambda\big(
    \rho_1\varepsilon_{\mathrm{leaf}}
    +\rho_2\varepsilon_{\mathrm{merge}}
    +\rho_3\varepsilon_{\mathrm{idemp}}\big).
\end{equation}
Setting $\rho_1=\rho_2=\rho_3=1$ recovers the main-text
$J_\lambda$; setting $\lambda=0$ gives pure oracle risk; setting
$\lambda=1$ gives the pure local-law projection term (for any fixed
share vector). The overall weight vector applied to the oracle term and
the three per-law terms is
$(1-\lambda,\,\lambda\rho_1,\,\lambda\rho_2,\,\lambda\rho_3)$, and the
objective bound flows through
\[
\text{Section 9 realization}
\;\to\;
\text{approximate local laws}
\;\to\;
\Delta_R \text{ / preference-gap bounds}.
\]
A fully-general per-budget weight vector $\boldsymbol{\lambda}=(\lambda_{\mathrm{leaf}},\lambda_{\mathrm{merge}},\lambda_{\mathrm{idemp}})$
that dispenses with the scalar/share decomposition subsumes the above
(with $\lambda_i := \lambda\rho_i$); we do not adopt it in the main text
because the scalar $\lambda\in[0,1]$ makes the iff endpoint theorems
cleaner to state.

This is the function-space reading of the training knob. Small
$\lambda$ prioritizes approximating the task oracle; large $\lambda$
prioritizes moving the realized operator into the local-law-constrained
subspace. The audit supplies the empirical certificate for where the
trained operator actually lands.

\subsection{Simplicity: Local-Law Weights Are a Projection}\label{app:projection-interpretation}

It is worth spelling out how simple the construction is when viewed
through the function-space lens. Fix a neural-operator class $\mathcal{N}$
--- the image of whatever parametric family is being trained --- and fix
an oracle $\fstar$ and tree $T$. The \emph{learnable local-law subspace}
is then just the intersection
\[
\mathrm{ExactLL}_{\mathcal{N}}(T,\fstar)
  \;:=\;
  \mathcal{N}\cap\mathrm{ExactLL}(T,\fstar),
\]
i.e.\ the operators that the class can represent \emph{and} that the
local laws permit. This is the same symbol already used in
Section~\ref{sec:theorems}. There is nothing in the training objective
that goes beyond pushing the learned operator into this set.

To make this precise, we separate two readings that are often conflated
in prose. The first --- ``local-law weights are a projection'' --- is
the set-level statement that the zero set of the local-law penalty is
exactly the exact local-law subspace. The second --- ``the
approximation error is structured by the local laws'' --- is the
pointwise statement that the only operators $g$ at which
$\mathrm{projectionPenalty}(g) = 0$ are the operators satisfying
C1/C2/C3. The following theorem records that these are the same
statement.

\begin{theorem}[Projection/structured-error equivalence]\label{thm:projection-iff}
Fix a tree $T$, an oracle $\fstar$, and a projection penalty
$\mathrm{projectionPenalty} : \Strings\rightsquigarrow\Strings\to\mathbb{R}_{\ge 0}$.
The following are equivalent:
\begin{enumerate}[leftmargin=1.5em,nosep]
\item \emph{(Projection framing.)}
$\{g : \mathrm{projectionPenalty}(g) = 0\} \;=\; \mathrm{ExactLL}(T,\fstar).$
\item \emph{(Structural-error framing.)} For every summarizer $g$,
$\mathrm{projectionPenalty}(g) = 0 \iff g \in \mathrm{ExactLL}(T,\fstar).$
\end{enumerate}
Moreover, for any fixed neural-operator class $\mathcal{N}$, both statements are
equivalent to
\[
\mathcal{N} \cap \{g : \mathrm{projectionPenalty}(g) = 0\}
  \;=\;
  \mathrm{ExactLL}_{\mathcal{N}}(T,\fstar).
\]
\end{theorem}

\noindent\emph{Proof.} Equivalence of (1) and (2) is set-extensionality
applied to the displayed sets. The class-restricted version follows by
intersecting both sides of (1) with $\mathcal{N}$ and using
$\mathcal{N}\cap\mathrm{ExactLL}(T,\fstar)=\mathrm{ExactLL}_{\mathcal{N}}(T,\fstar)$.
\qed

\medskip
In words: the local-law weights we add to the loss are, up to
faithfulness of the penalty, \emph{exactly} a projection operator onto
the learnable local-law subspace. The $\lambda\to 1$ limit of
$J_\lambda$ behaves like an indicator penalty for that set. Large
$\lambda$ is hard projection; small $\lambda$ is soft projection traded
off against oracle fit. The method is not an arbitrary regularizer; it
is a weighted penalty for leaving the subspace the theorems actually
need.

The operational $J_\lambda$ that the training loop uses is the
finite-sample realization of this picture with the \emph{approximate}
local-law bundle in place of an exact indicator. The audit of
\S\ref{sec:estimation} is what closes the loop: it reports where in
$\mathrm{ApproxLL}(T,\fstar;\varepsilon)$ the trained $g$ has landed,
and the bound in \S\ref{sec:discussion} converts that audited
$\varepsilon$ into the IPW term of Equation~\eqref{eq:discussion-bound}.

\paragraph{Projection framing $\iff$ structural assumption on the approximation error.}
Theorem~\ref{thm:projection-iff} is the precise content of the
equivalence the main text refers to as the projection/structured-error
iff. When we introduce $\lambda$-weighted local-law terms in training,
we are making a statement about what the trained operator is
approximating. Calling it a ``projection'' makes the \emph{target} (the
learnable local-law subspace) visible; calling it ``an assumption about
the structure of the approximation error'' makes the \emph{content of
the assumption} visible. They are the same assumption, and the iff is
what lets us move between them without sleight of hand.
